% ========================================================
%  [课程名]  [考试类型]（含参考答案）
%  覆盖范围：[知识点1]、[知识点2]、[知识点3]……
% ========================================================
\documentclass[12pt, a4paper]{article}

% ---- 中文支持 ----
\usepackage[UTF8]{ctex}

% ---- 数学 ----
\usepackage{amsmath, amssymb, amsthm}
\usepackage{mathtools}

% ---- 页面布局 ----
\usepackage[top=2.5cm, bottom=2.5cm, left=2.8cm, right=2.8cm]{geometry}

% ---- 表格 ----
\usepackage{booktabs, array, multirow, makecell}
\usepackage{tabularx}

% ---- 页眉页脚 ----
\usepackage{fancyhdr}
\usepackage{lastpage}
\pagestyle{fancy}
\fancyhf{}
\renewcommand{\headrulewidth}{0.6pt}
\lhead{\small [课程名]·[考试类型]}
\rhead{\small 共 \pageref{LastPage} 页}
\cfoot{\small 第 \thepage\ 页}

% ---- 计数器与样式 ----
\usepackage{enumitem}
\usepackage{xcolor}
\usepackage{tcolorbox}
\tcbuselibrary{skins, breakable}

% ---- 填空线 ----
\newcommand{\blank}[1]{\underline{\hspace{#1}}}

% ---- 答案盒子 ----
\newtcolorbox{answerbox}[1][]{
  breakable,
  colback=gray!6,
  colframe=gray!50,
  fonttitle=\bfseries\small,
  title=参考答案,
  #1
}

% ---- 题号格式 ----
\newcounter{bigq}
\newcommand{\BigQ}[1]{%
  \stepcounter{bigq}%
  \vspace{6pt}%
  \noindent{\large\bfseries 第\chinese{bigq}大题\quad #1}%
  \vspace{4pt}\par
}

\newcounter{subq}[bigq]
\newcommand{\SubQ}{%
  \stepcounter{subq}%
  \noindent\textbf{（\arabic{subq}）}\quad
}

% 分隔线
\newcommand{\examrule}{\noindent\rule{\linewidth}{0.6pt}}

% ---- 文档开始 ----
\begin{document}

% ============================================================
%  封面
% ============================================================
\begin{center}
  {\LARGE\bfseries [课程名]}\\[8pt]
  {\large\bfseries [考试类型]}\\[16pt]
  \begin{tabular}{lll}
    考试时间：[时长] 分钟 & \quad & 满分：[分数] 分 \\[4pt]
    \multicolumn{3}{l}{注意事项：请将答案写在答题纸指定区域，写在本卷上无效。}
  \end{tabular}
  \vspace{4pt}

  \examrule

  \vspace{6pt}
  \begin{tabular}{|c|c|c|c|c|c|c|c|c|}
    \hline
    \textbf{题号} & 一 & 二 & 三 & 四 & 五 & 六 & 七 & \textbf{总分} \\
    \hline
    \textbf{满分} & & & & & & & & \\
    \hline
    \textbf{得分} & & & & & & & & \\
    \hline
  \end{tabular}
\end{center}

\vspace{10pt}
\examrule
\vspace{6pt}

% ============================================================
%  第一大题
% ============================================================
\BigQ{[题型名称]（每题 x 分，共 x 分）}

\SubQ [题目内容]

\vspace{6pt}
\SubQ [题目内容]

\vspace{10pt}
\examrule

% ============================================================
%  第二大题
% ============================================================
\BigQ{[题型名称]（共 x 分）}

\SubQ（x 分）[题目内容]

\vspace{10pt}
\examrule

% ============================================================
%  第三大题
% ============================================================
\BigQ{[题型名称]（共 x 分）}

\SubQ（x 分）[题目内容]

\vspace{10pt}
\examrule

% ============================================================
%  更多大题……
% ============================================================
% 复制以下模板添加新大题：
%
% \BigQ{[题型名称]（共 x 分）}
%
% \SubQ（x 分）[题目内容]
%
% \vspace{10pt}
% \examrule

\vspace{16pt}
\examrule
\vspace{4pt}
\begin{center}
  {\large\bfseries ——试题到此结束，请认真检查——}
\end{center}
\vspace{4pt}
\examrule

% ============================================================
%  参考答案
% ============================================================
\newpage
\setcounter{bigq}{0}

\begin{center}
  {\LARGE\bfseries 参考答案与评分标准}
\end{center}
\vspace{8pt}
\examrule

% -------- 第一大题参考答案 --------
\BigQ{[题型名称]参考答案}

\begin{answerbox}

[答案内容]

\end{answerbox}

% -------- 第二大题参考答案 --------
\BigQ{[题型名称]参考答案}

\begin{answerbox}

[答案内容]

\end{answerbox}

% -------- 第三大题参考答案 --------
\BigQ{[题型名称]参考答案}

\begin{answerbox}

[答案内容]

\end{answerbox}

% 更多答案……复制以下模板：
%
% \BigQ{[题型名称]参考答案}
%
% \begin{answerbox}
%
% [答案内容]
%
% \end{answerbox}

\vspace{12pt}
\examrule
\begin{center}
  \textit{参考答案仅供参考，评分时应酌情给分，步骤正确可得过程分。}
\end{center}

\end{document}
